% ###################
% ## Filename: sCM_mode.tex
% ## Location: D:\Project\Prej_Teaching\ISE\SupplyChainMgt\presentation\scm_06_mode
% ## Created: Jan 26, 2011
% ## Last updated: Feb 2, 2010
% ## Figure Location: ..\Figures\
% ## Note: Presentation, Teaching
% ###################

\documentclass[10pt,presentation]{beamer}
%\documentclass[10pt,handout]{beamer}
%\documentclass[10pt,handout,notes=show]{beamer}


%==============================================================================================================
\usepackage{fontspec}
\usepackage{xunicode}
\usepackage{xltxtra}
\usepackage{array}
\usepackage{xcolor}
\usepackage{verbatim}
\usepackage{beamerthemesplit}
\usepackage{beamerbasetemplates}
\usepackage{graphics}
\usepackage{graphicx}
\usepackage{hyperref}
\usepackage{multirow}
\usepackage{tikz,pgfplots}
\usepackage{algorithm}
\usepackage{algorithmic}
\usepackage{soul,xcolor}
\usepackage{colortbl}
\usepackage{multicol}

\XeTeXlinebreaklocale"th"
\defaultfontfeatures{Scale=1.5, Mapping=tex-text} %{Scale=1.23}



\usetheme{Madrid}  % JuanLesPins  Rochester
           %  Berkeley  Palo Alto    with sidebar and top
           % Goettingen  Marburg     with sidebar
           %  Copenhagen  Luebeck  Warsaw

\setbeamertemplate{headline}
{%
    \begin{beamercolorbox}{section in head/foot}
        \vskip-2pt \hskip5pt
        \insertsectionnavigationhorizontal{\textwidth}{}{}
        \hskip5pt \vskip-2pt
    \end{beamercolorbox}
}%
\setbeamertemplate{footline}
{%
    \begin{beamercolorbox}{section in head/foot}
        \vskip-2pt \hskip5pt
        \insertshorttitle \hfill \insertshortdate \hskip10pt \insertframenumber / \inserttotalframenumber
        \vskip-2pt\hskip5pt
    \end{beamercolorbox}
}%
\definecolor{alertedOrange}{RGB}{255,36,0}
\setbeamercolor{alerted text}{fg=alertedOrange}
\useinnertheme{rounded}
\usecolortheme{rose}
%\setbeamercovered{transparent}
\usefonttheme{structuresmallcapsserif}
\definecolor{LRed}{rgb} {1,.8,.8}
\definecolor{LBlue}{rgb}{.8,.8,1}

%%%%%%%%%%%%%%%%%%%%%%%%%%%%%%%%%%%%%%%%%%%%%%%%%%%%%%%%%%%%%%%%%%%%%%%%%%%%%%%
% Notation:
\newtheorem{observation}{Observation}
\newcommand{\figureLocation}[1]{../figures/#1}
\newcommand{\figureLoc}[1]{../figures/#1}
\newcommand{\refLocation}[1]{../../../../Reference/#1}
\newcommand{\NextLine}{\vskip 0.05\textheight}
\newcommand{\NextHalfLine}{\vskip 0.02\textheight}
\newcommand{\RemoveLine}{\vskip -0.1\textheight}
\newcommand{\RemoveHalfLine}{\vskip -0.02\textheight}
\newcommand{\BackTab}{\hskip -0.15\textwidth}
\newcommand{\HalfBackTab}{\hskip -0.11\textwidth}
\newcommand{\DualFigureSize}{1.10\textwidth}
\newcommand{\MiniPageSize}{0.52\textwidth}
\newcommand{\code}[1]{\textcolor{blue}{#1} }
\newcommand{\rightNote}[2]{\hfill \uncover<#1>{\textbf{\code{\#\# {#2}  }}}}
\newcolumntype{P}[1]{>{\raggedright\arraybackslash}p{#1}<{}}

\author[Kittithreerapronchai]{%
  Oran~Kittithreerapronchai\inst{1} }
\institute[IE Chula]{
  \inst{1}%
  Department of Industrial Engineering, Chulalongkorn University\\
  Bangkok 10330 THAILAND \\
  }
\subject{}
\date{}

\setbeamertemplate{enumerate items}[default]
\setbeamertemplate{bibliography item}[text]
\bibliographystyle{alpha}
\newcommand{\textFont}[1]{\footnotesize{#1}}
\setcounter{MaxMatrixCols}{20}

%%%--------------------------------

\makeatletter
\defbeamertemplate{note page}{infolines}
{%
  {%
    \small
    \insertvrule{.25\paperheight}{white!90!black}
    \vskip-.25\paperheight
    \nointerlineskip
    \vbox{
      \hfill\insertslideintonotes{0.25}\hskip-\Gm@rmargin\hskip0pt%
      \vskip-0.25\paperheight%
      \nointerlineskip
      \begin{pgfpicture}{0cm}{0cm}{0cm}{0cm}
        \begin{pgflowlevelscope}{\pgftransformrotate{90}}
          {\pgftransformshift{\pgfpoint{-2cm}{0.2cm}}%
          \pgftext[base,left]{\footnotesize\the\year-\ifnum\month<10\relax0\fi\the\month-\ifnum\day<10\relax0\fi\the\day}}
        \end{pgflowlevelscope}
      \end{pgfpicture}}
    \nointerlineskip
    \vbox to .25\paperheight{\vskip0.5em
      \hbox{\insertshorttitle[width=8cm]}%
      \setbox\beamer@tempbox=\hbox{\insertsection}%
      \hbox{\ifdim\wd\beamer@tempbox>1pt{\hskip4pt\raise3pt\hbox{\vrule
            width0.4pt height7pt\vrule width 9pt
            height0.4pt}}\hskip1pt\hbox{\begin{minipage}[t]{7.5cm}\def\breakhere{}\insertsection\end{minipage}}\fi%
      }%
      \setbox\beamer@tempbox=\hbox{\insertsubsection}%
      \hbox{\ifdim\wd\beamer@tempbox>1pt{\hskip17.4pt\raise3pt\hbox{\vrule
            width0.4pt height7pt\vrule width 9pt
            height0.4pt}}\hskip1pt\hbox{\begin{minipage}[t]{7.5cm}\def\breakhere{}\insertsubsection\end{minipage}}\fi%
      }%
      \setbox\beamer@tempbox=\hbox{\insertshortframetitle}%
      \hbox{\ifdim\wd\beamer@tempbox>1pt{\hskip30.8pt\raise3pt\hbox{\vrule
            width0.4pt height7pt\vrule width 9pt
            height0.4pt}}\hskip1pt\hbox{\insertshortframetitle[width=7cm]}\fi%
      }%
      \vfil}%
  }%
  \vskip.25em
  \nointerlineskip
  %\setmainfont{TH Sarabun New}
  \begin{minipage}{\textwidth} % this is an addition
  \footnotesize{\fontspec{TH Sarabun New} \insertnote}
  \end{minipage}               % this is an addition
  %\setmainfont{Computer Modern}

}
\makeatother

\setbeamertemplate{note page}[infolines]

\title[LSCM v3.0]{Lecture 07: \\ Modes of Transportation  \\ ~~}
\hypersetup{
  pdftitle={07 Modes of Transportation},
  pdfsubject={2104509lscm & 2184408scm handout},
  pdfproducer={XeLatex/WinEdt},
  pdfcreator={oran.k@chula.ac.th}
}


\begin{document}

%==============================================================================================================
\begin{frame}[t]
  \titlepage \label{Slide:Title}

  \NextLine   \NextLine
  \begin{flushright}
    \tiny{ \textbf{last updated:} \today } \hskip3ex
  \end{flushright}

  \note{
    \noindent\textbf{Purpose:} Overview and Introduction to modes of transportation\par
    \noindent\textbf{Leading Question:} Where have you encountered this kind of supply-chain decision?\par
    \noindent\textbf{Teaching cue:} Get one example and identify who owns the decision and what outcome they want.\par
    \begin{itemize}
      \item \textbf{Goal1:} understand each mode/ requirement/ situation
      \item \textbf{Goal2:} relationship between modes and product|country
      \item \textbf{Goal3:} critical Thailand infrastructure of each mode
    \end{itemize}
  }
\end{frame}



%==============================================================================================================
\begin{frame}[t]
  \frametitle{Outline}

  \tableofcontents
  \vfill
  \begin{flushright}
    \tiny{ \textbf{Key Ref.:}~\cite{jacobs2010operations}~\cite{ballou2007evolution} \cite{chopra2007supply} \cite{goetschalckx2011supply}} \hskip3ex
  \end{flushright}

  \note{
    \noindent\textbf{Purpose:} Overview of Transportation\par
    \noindent\textbf{Leading Question:} Which topic would you consult first for the decision raised at the opening?\par
    \noindent\textbf{Teaching cue:} Point to the relevant section and ask which question should be resolved by the end.\par
    \begin{itemize}
      \item \textbf{Statistic} introduction to standard modes and assess your knowledge in each modes
      \item \textbf{Issue:} detail operation and critical components/nature of each mode
      \item \textbf{Infrastructure:} a case study of Thappline
    \end{itemize}
  }
\end{frame}

%==============================================================================================================
\section[Mode 101]{Statistics of Modes}
%==============================================================================================================

\begin{frame}
  \frametitle{US Domestic Statistic 2013}
  \small{
  \begin{tabular}{lrrr|rrr} \hline		
    Mode               & Tons      & Value    & Ton-miles$^{07}$ & \multicolumn{3}{|c}{Percentage} \\
                       & (million) & (billion \$)& (billion) & Tons  & Value  &Ton-miles$^{07}$ \\ \hline
    Truck	             &13,732	 &10,841	&2349	&77\% &75\% 	&41\% \\
    Rail	             &1,681	   &424	    &1522	&9\% 	&3\% 	&27\% \\
    Water	             &410	     &131	    &450	&2\% 	&1\% 	&8\% \\
    Air	               &3	       &134	    &9	  &0\% 	&1\% 	&0\%  \\
    Mail               &459 &1,695	  &469	&3\% 	&12\% 	&8\% \\
    Pipeline	         &1,391	   &1,003	  &855	&8\% 	&7\% 	&15\% \\
    Other              &274	     &270	    &86	  &2\% 	&2\% 	&1\% \\ \hline
  \end{tabular}
  }

%Truck	11,520,318	12,421,511
%Rail	1,738,346	690,459
%Water	766,322	363,500
%Air (including truck-air)	5,871	591,253
%Multiple modes and mail	495,680	2,328,112
%Pipeline	3,049,857	942,007
%Other and unknown	39,210	97,633


  \begin{flushright}
    \tiny{\textbf{source:} US Bureau of Transportation}
  \end{flushright}

  \note{
    \noindent\textbf{Purpose:} Overview of modes in US and how to measure\par
    \noindent\textbf{Leading Question:} Why can payload, shipment value, and ton-miles give different mode shares?\par
    \noindent\textbf{Teaching cue:} Compare the denominators and the table's year labels before drawing a market conclusion.\par
    \begin{itemize}
      \item \textbf{Truck:} major in all terms, esp value $\rightarrow$ why? (as last mile)
      \item \textbf{Rail:} low value \& bulky products (grain/car/oil), long haul $\rightarrow$ unique to USA
      \item \textbf{Water:} two agriculture waterways: Mississippi System and Great lakes System $\rightarrow$ which items
      \item \textbf{Air:} expensive mode for parcels and special items $\rightarrow$ do applicable to Thailand?
      \item \textbf{Pipeline:} unique infrastructure for energy and utility (oil/gas/coal/water)
      \item \textbf{Parcel/Intramodel/Mail:}
    \end{itemize}
  }
  \end{frame}


%%==============================================================================================================
\begin{frame}
  \frametitle{US freight map}

  \begin{center}
    \includegraphics[width=0.95\textwidth]{\figureLocation{scm_mode/US_map_freigt_mode.png}}
  \end{center}

  \begin{flushright}
    \tiny{\textbf{source:} US Bureau of Transportation}
  \end{flushright}

  \note{
    \noindent\textbf{Purpose:} Detail features by distance\par
    \noindent\textbf{Leading Question:} Which shipment movements are emphasized on the freight map?\par
    \noindent\textbf{Teaching cue:} Identify mode, geography, and year from the original legend before comparing corridors.\par
    \begin{itemize}
      \item \textbf{Distance:} long $\rightarrow$ cheap (exception of Air)
      \item \textbf{Value:} truck and pipeline
      \item \textbf{Tonage:}
      \item \textbf{Other:} OD, nature of goods, time,
    \end{itemize}
    \begin{itemize}
      \item \textbf{case:} similar to weight; low value
      \item \textbf{points:} not just truck, depends on mode of transportation; different products,
    \end{itemize}
  }
  \end{frame}

%%==============================================================================================================
%\begin{frame}
%  \frametitle{US freight by distance band \fo}
%
%  \begin{tabular}{cc}
%    \includegraphics [width=0.45\textwidth]{\figureLocation{scm_mode/valueBand2007.png}}
%    &
%    \includegraphics [width=0.45\textwidth]{\figureLocation{scm_mode/tonBand2007.png}}
%    \\
%    \includegraphics [width=0.45\textwidth]{\figureLocation{scm_mode/tonMileBand2007.png}}
%    &
%    \includegraphics [width=0.45\textwidth]{\figureLocation{scm_mode/legendBand2007.png}}
%  \end{tabular}
%
%
%  \begin{flushright}
%    \tiny{\textbf{source:} US Bureau of Transportation}
%  \end{flushright}
%
%  \note{
%    \textbf{Purpose:} Detail features by distance \\
%    \textbf{Leading Question:} how you choose modes of transportation for your goods?
%    \begin{itemize}
%      \item \textbf{Distance:} long $\rightarrow$ cheap (exception of Air)
%      \item \textbf{Value:} truck and pipeline
%      \item \textbf{Tonage:}
%      \item \textbf{Other:} OD, nature of goods, time,
%    \end{itemize}
%    \begin{itemize}
%      \item \textbf{case:} similar to weight; low value
%      \item \textbf{points:} not just truck, depends on mode of transportation; different products,
%    \end{itemize}
%  }
%  \end{frame}

%==============================================================================================================
\begin{frame}
  \frametitle{How about Thailand?}

  \begin{center}\small{
  \begin{tabular}{lrrcc} \cline{1-3}
    Mode  & USA     & Thailand & ~ & \multirow{6}{*}{ \begin{minipage}{0.40\textwidth} \includegraphics [width=\textwidth]{\figureLocation{scm_mode/modeUSA_THAI.png}}  \end{minipage}}\\ \cline{1-3}
    Truck	& 13,732	& 426.1 & \\
    Rail	& 1,681	  & 11.8 \\
    Water	& 410	    & 91.7 \\
    Air	  & 3	      & 0.1 \\
    Pipeline	& 1,391	& - \\
    Other	& 733	    & - \\ \cline{1-3}
  \end{tabular} }
  \end{center}
        \begin{flushright}
          \tiny{\textbf{source:} US Bureau of Transportation and TH Ministry of Transportation}
        \end{flushright}

  \begin{block}{What do these data imply?}{}{}
    \begin{itemize}
      \item $\exists$ relationship between mode and type of product
      \item Compare reported shares using the same units, year, and mode coverage
      \item The table compares national freight volumes; last-mile and trade shares need separate data
    \end{itemize}
  \end{block}

  \note{
    \noindent\textbf{Purpose:} Introduce mode nature in Thailand\par
    \noindent\textbf{Leading Question:} What normalization is needed to compare freight mode between countries?\par
    \noindent\textbf{Teaching cue:} Check units, dates, modes covered, and treatment of missing values before using the figure.\par
    \begin{itemize}
      \item \textbf{Truck:} has the largest volume among the modes reported for both countries
      \item \textbf{USA:} uses more pipeline and rail
      \item \textbf{Thailand:} uses more water
    \end{itemize}
    \textsc{\textbf{Note}}
    \begin{itemize}
      \item \textbf{Thailand:} Truck 80\%; river 9\%; sea 8.5\%; rail 2.5\%
      \item \textbf{points:} not just truck, depends on mode of transportation; different products,
    \end{itemize}
  }
  \end{frame}

%==============================================================================================================
\begin{frame}
  \frametitle{What is a suitable mode?}

  \begin{itemize}
    \item \textbf{Air:} expensive, any where, quick, safely (e.g., CPU, produces)
    \item \textbf{Rail:} \alert{city-to-city}, bulky (e.g., ores, coals, grains, automobiles)
    \item \textbf{Truck:} accessability, \alert{control}, \alert{responsive}
    \item \textbf{Water:} international, bulky, cheap, \alert{low emission} (e.g., cargo)
    \item \textbf{Pipeline:} special infrastructure (e.g., oil, natural gas)
    \item \textbf{Parcel:} postal services (e.g., Dell computer, Amazon)
    \item \textbf{Intermodal:} more than one mode (e.g., rail+truck, water+rail)
  \end{itemize}

  \NextLine
  \textbf{Modern Marvel: Delivery it}
  \begin{itemize}
    \item \textbf{Story:} Extreme delivery and transportation
    \item \textbf{Observation:}
    \begin{itemize}
      \item What is information necessary for planning and execute each type of delivery?
      \item What are equipments and/or innovation in each type of delivery?
      \item What are challenges in each type of delivery?
    \end{itemize}
  \end{itemize}

  \note{
    \noindent\textbf{Purpose:} Overview/Summary of each mode\par
    \noindent\textbf{Leading Question:} Which delivery promise and commodity constraints rule out a mode?\par
    \noindent\textbf{Teaching cue:} Filter by feasibility first, then compare door-to-door cost and reliability.\par
    \textbf{Trivial of each mode}
      \begin{itemize}
        \item \textbf{Air} charter flight, mixed passenger+cargo flight
        \item \textbf{Rail} common unitload for rail and truck $\rightarrow$ carload (CL)
        \item \textbf{Truck} ownership for vehicle, size of truck, pricing,  TL-VS-LTL; Class
        \item \textbf{Water} 80\% of international freight, slowest
        \item \textbf{Pipeline} grounded coal
    \end{itemize}
  }

  \end{frame}

%==============================================================================================================
\begin{frame}
  \frametitle{Air System: High speed and expensive }

  \begin{itemize}
    \item \textbf{Historical:} recent last couple of decades
    \item \textbf{Concurrent Uses:} passenger, air mail, produces, electronic gadgets, jewelry
    \item \textbf{Opportunity:} \alert{passenge+cargo}, revenue management
    \item \textbf{Issues:} \alert{security}, hub location, crew scheduling, fuel consumption, maintenance
    \item \textbf{Carriers:} UPS, FedEx, TNT, carter flight
  \end{itemize}

  \note{
    \noindent\textbf{Purpose:} Overview of Air transportation\par
    \noindent\textbf{Leading Question:} When can the added cost of air freight reduce total supply-chain cost?\par
    \noindent\textbf{Teaching cue:} Compare transit time, inventory exposure, and cost of a missed delivery.\par
    \begin{itemize}
      \item \textbf{Airport Free Zone:} no tax, yet. for import-export + avoid tax refund redtap
      \item \textbf{THCargo:} THAI , charge everything, always look at manual for answer
      \item \textbf{Cargo:} produce , gold, harzard chem, bicycle
      \item \textbf{authority of caption:} own nation, priority between goods VS passenger, similarity between air and sea
    \end{itemize}
  }
  \end{frame}

%==============================================================================================================
\begin{frame}
  \frametitle{Main Air Shipping Routes}

  \begin{center}
    \includegraphics[width=0.80\textwidth]{\figureLocation{scm_mode/mainAirRoutes.png}}
  \end{center}

  \note{
    \noindent\textbf{Purpose:} Air transportation system\par
    \noindent\textbf{Leading Question:} What opportunities and constraints follow from an airport hub?\par
    \noindent\textbf{Teaching cue:} Compare frequency, connections, ground transfer, and available cargo capacity.\par
    \begin{itemize}
      \item \textbf{dominate by:} USA-EUR transportation
      \item \textbf{jet stream:} JPN to USA (during WWII)
      \item \textbf{3D lane of airplane:} depend on height
    \end{itemize}
  }
  \end{frame}
%==============================================================================================================
\begin{frame}
  \frametitle{Rail System: Opportunity or Legacy}

  \begin{itemize}
    \item \textbf{Historical:} key infrastructure, creating many metropolis
    \item \textbf{Pricing:} High fixed cost (locomotive, rail way)
    \item \textbf{Concurrent Uses:} vary depends largely on region
    \begin{itemize}
      \item \textbf{Europe:} high-speed passengers
      \item \textbf{N. America:} bulky freight + passengers
      \item \textbf{Asia Pacific:} budget transportation for passengers
    \end{itemize}
    \item \textbf{Opportunity:} \alert{intermodal transportation}, landbridge
    \item \textbf{Issues:} locomotive management, inventory, \alert{scheduling}, gauge
    \item \textbf{Carriers:} Norfolk Southern, CSX Transportation
  \end{itemize}

  \note{
    \noindent\textbf{Purpose:} Overview of Rail System\par
    \noindent\textbf{Leading Question:} Which shipment characteristics make rail attractive?\par
    \noindent\textbf{Teaching cue:} Compare density and distance with transloading cost and time.\par
    \begin{itemize}
      \item \textbf{Note:} 60\% transportation cost is fuel
      \item \textbf{rail:} underestimate mode of transportation; obsolete
      \item \textbf{benefits:} connect one metropolis to other; scalable capacity
      \item \textbf{rail size:} CL(carload), no design, no hub
      \item Very small number of carriers; likely only one will be able to serve any one customer location
      \item Trend is to merge smaller companies into larger ones with ultimate goal of having perhaps two transcontinental rail carriers
    \end{itemize}
  }
  \end{frame}

%==============================================================================================================
\begin{frame}
  \frametitle{Track gauge and Double Track}

  \begin{center}
    \only<1|handout:1>{  \includegraphics[width=\textwidth]{\figureLocation{scm_mode/trackGauge.png}}}
    \only<3-|handout:0>{ \includegraphics[width=0.8\textwidth]{\figureLocation{scm_mode/doubleRailThailand.jpg}}}
    \only<2|handout:0>{ \includegraphics[width=0.8\textwidth]{\figureLocation{scm_mode/interModal.png}}}
  \end{center}

  \note{
    \noindent\textbf{Purpose:} Thailand rail infrastructure\par
    \noindent\textbf{Leading Question:} What operational issue can arise where track gauges differ?\par
    \noindent\textbf{Teaching cue:} Trace a handoff and ask how transfer time and capacity alter the route plan.\par
    \begin{itemize}
      \item \textbf{Infrastructure problem:} how go south, Inland Container Depot
      \item \textbf{Railway Cost:} 100.0 M.THB /km (dual track)
      \item \textbf{Thailand rail length} 4,070 km; metre gauge (1.0 m) same as every else
    \end{itemize}
  }
  \end{frame}

%==============================================================================================================
\begin{frame}
  \frametitle{Types of Rail Car }

  \begin{tabular}{cc}
    \includegraphics [width=0.45\textwidth]{\figureLocation{scm_mode/flatCar.jpg}}
    &
    \includegraphics [width=0.45\textwidth]{\figureLocation{scm_mode/boxCar.jpg}}
    \\
    \includegraphics [width=0.45\textwidth]{\figureLocation{scm_mode/gondolaCar.jpg}}
    &
    \includegraphics [width=0.45\textwidth]{\figureLocation{scm_mode/hopperCar.jpeg}}
  \end{tabular}


  \note{
    \noindent\textbf{Purpose:} Type of cargo\par
    \noindent\textbf{Leading Question:} How does product handling constrain rail-car selection?\par
    \noindent\textbf{Teaching cue:} Match one commodity to the equipment and state a loading or unloading constraint.\par
    \begin{itemize}
      \item \textbf{Bogie Type:}
      \item \textbf{Daqin Railway, China} 130 tonnes of coal per wagon and tens of thousands of tonnes per train.
    \end{itemize}
  }
  \end{frame}

%==============================================================================================================
\begin{frame}
  \frametitle{Truck System: Dominated Mode}

  \begin{itemize}
    \item \textbf{Classification:} less-than truckload (LTL), Full truckload (TL)
    \item \textbf{Pricing:} weight, weight-distance, area
    \begin{itemize}
      \item \textbf{weight:} TL \textbf{VS} LTL
      \item \textbf{distance:} national \textbf{VS} regional
      \item \textbf{class:} hazard, special equipment
    \end{itemize}
    \item \textbf{Opportunity:} \alert{network design}, consolidation, realtime management
    \item \textbf{Issues:} \alert{fleet management}, gas price, backhaul
    \item \textbf{Carriers:} Schneider National, JB Hunt, Ryder Integrated
  \end{itemize}

  \note{
    \noindent\textbf{Purpose:} Overview of truck system\par
    \noindent\textbf{Leading Question:} Which delivery conditions make truck service attractive despite unit costs?\par
    \noindent\textbf{Teaching cue:} Compare door access, flexibility, shipment size, and route density.\par
    \begin{itemize}
      \item High accessibility
      \item Cost of construction: 20 M.THB /km
      \item weight limit 50 ton for two trailer; 15 ton for 5W; 25ton for 10w
      \item TH Road 1 (1000) + Road 2 (510) + Road 3 (488) + Road 4 (1222)
    \end{itemize}
  }
  \end{frame}

%==============================================================================================================
\begin{frame}
  \frametitle{Water System: international trade }

  \begin{itemize}
    \item \textbf{Important:} dominated international trade
    \item \textbf{Pricing:} volume, class, term (e.g., FOB)
    \item \textbf{Opportunity:} \alert{efficiency}, environmental friendly, safety
    \item \textbf{Issues:} \alert{fleet management}, pirate, capacity, leadtime
    \item \textbf{Carriers:} Maersk Sealand, Evergreen, Hanjin Shipping Co
    \item \textbf{nautical VS statute mile}  1.0 nautical mile/hr = 1.0 knot = 1.8 km/hr
  \end{itemize}


  \note{
    \noindent\textbf{Purpose:} Overview of water system\par
    \noindent\textbf{Leading Question:} How can water transport change inventory and service requirements?\par
    \noindent\textbf{Teaching cue:} Compare sailing frequency and port handling with linehaul cost and inland connections.\par
    \begin{itemize}
      \item charge for everything similar to air (extra waiting, dock area, equipment)
      \item costal and river; River: 200/3,701 km
    \end{itemize}
  }
  \end{frame}

%==============================================================================================================
\begin{frame}
  \frametitle{Main Maritime Shipping Routes}

  \begin{center}
    \includegraphics[width=0.80\textwidth]{\figureLocation{scm_mode/mainMaritimeRoutes.png}}
  \end{center}

  \note{
    \noindent\textbf{Purpose:} infrastructure of water system\par
    \noindent\textbf{Leading Question:} How does a maritime service route affect total transit time?\par
    \noindent\textbf{Teaching cue:} Trace a port-to-port leg, possible transshipment, and inland legs.\par
    \begin{itemize}
      \item \textbf{key transit point:} strait of Hormuz, Malaccar (Singapore), gibraltar. CapeTown, Bosporus
      \item \textbf{Priate:} Nigeria, Bangladesh, Malagar, Aden
      \item \textbf{port cap}
    \end{itemize}
  }
  \end{frame}

%==============================================================================================================
\begin{frame}
  \frametitle{General type of marine vessels \forIE}

  \begin{tabular}{ccc}
    \includegraphics [width=0.30\textwidth]{\figureLocation{scm_mode/riverVessel.jpg}}
    &
    \includegraphics [width=0.30\textwidth]{\figureLocation{scm_mode/costalVessel.jpg}}
    &
    \includegraphics [width=0.30\textwidth]{\figureLocation{scm_mode/oceanVessel.jpg}}    \\
    \multicolumn{3}{c}{
      \only<1|handout:1>{\includegraphics [width=0.60\textwidth]{\figureLocation{scm_mode/marineType01.png}}}
      \only<2|handout:0>{\includegraphics [width=0.60\textwidth]{\figureLocation{scm_mode/marineType02.png}}}
      \only<3-|handout:0>{\includegraphics [width=0.60\textwidth]{\figureLocation{scm_mode/dryBulkCarrierSize.jpg}}}
      }

  \end{tabular}

  \begin{flushright}
    \tiny{\textbf{source:} shippipedia.com}
  \end{flushright}

  \note{
    \noindent\textbf{Purpose:} maritime vessel\par
    \noindent\textbf{Leading Question:} Which vessel feature must match the freight and port?\par
    \noindent\textbf{Teaching cue:} Name the freight unit, handling method, and access requirement before choosing a vessel type.\par
    \begin{itemize}
      \item \textbf{Depth:} bigger sea + stronger wave, bigger depth + larger capacity $\rightarrow$ diff ocean, costal, and river
      \item \textbf{Name:} vessel type $\rightarrow$ infrastructure
    \end{itemize}
  }
  \end{frame}

%==============================================================================================================
\begin{frame}
  \frametitle{Type of vessels by freight \forIE}

  \begin{center} \footnotesize{
  \begin{tabular}{lrlrlr}
    \multicolumn{2}{c}{\includegraphics [width=0.27\textwidth]{\figureLocation{scm_mode/tankerShip.jpg}}}
    &
    \multicolumn{2}{c}{\includegraphics [width=0.27\textwidth]{\figureLocation{scm_mode/bulkerShip.jpg}} }
    &
    \multicolumn{2}{c}{\includegraphics [width=0.27\textwidth]{\figureLocation{scm_mode/containerShip.jpg}} }\\
    ~ \\
    Tanker	  & DWT	            & Bulker	  & DWT	          & Container 	& TEU	\\ \hline
    Seawaymax	& 10.0k–60.0k	  & Handysize	& 10.0k–35.0k	& Panamax	    & 4.5k–5.0k	\\
    Panamax	  & 60.0k–80.0k	  & Handymax	& 35.0k–59.0k	& New Panamax & 5.0k–10.0k	\\
    Aframax	  & 80.0k–120.0k	& Panamax	  & 60.0k–80.0k	& Suezmax	    & 10.0k-15.5k	\\
    Suezmax	  & 120.0k–200.0k	& Capesize	& 80.k +	      & 	& \\ \hline
    Speed     & 10-15 knots     &           & 12-15 knots   &               & 15-25 knots \\ \hline
  \end{tabular} }
  \end{center}

  \begin{flushright}
    \tiny{\textbf{source:} shippipedia.com}
  \end{flushright}

  \note{
    \noindent\textbf{Purpose:} Overview Inventory in SCM\par
    \noindent\textbf{Leading Question:} How does a cargo property change loading or storage needs?\par
    \noindent\textbf{Teaching cue:} Match the product's physical handling requirements to a suitable ship category.\par
    \begin{itemize}
      \item \textbf{Seawaymax} can pass lock within N.America (CAN-USA) at Great Lake
      \item \textbf{Aframax} 	Tankers up to 120 000 deadweight tonnes, designed for the entry into North American Atlantic ports, with a maximum load. Afra = American Freight Rate Association.
      \item \textbf{Panamax} 	Vessels that meet at least one of the following criteria, width 106 ft (30.5 m), length 965 feet (294.1 meters) deep or 39.5 feet (12.04 meters).Panamax is also a term for the bulk of 50000-80000 tons dw. Approximately 80,000 tons dw is the maximum size for a fully loaded vessel to cross the Panama Canal, hence the name “panmax”
      \item \textbf{Suezmax} 	Tankers of about 150,000 tons dw. Approximately 150,000 tons dw is the maximum size for a fully loaded vessel to cross the Suez Canal, hence the name “Suezmax”.
      \item \textbf{Capesize} 	Bulk of the order of 100 000-150 000 tons deadweight. cannot pass suze and panama
      \item \textbf{Handysize} 	Bulk of 20000-50000 tonnes deadweight.
    \end{itemize}
  }
  \end{frame}

%==============================================================================================================
\begin{frame}
  \frametitle{Main seaports in Thailand}

  \begin{center} \footnotesize{
  \begin{tabular}{cc}
    \includegraphics [width=0.6\textwidth]{\figureLocation{scm_mode/thaiSeaports.png}}
    &
    \includegraphics [width=0.3\textwidth]{\figureLocation{scm_mode/thaiSeaportsMap.png}}
  \end{tabular} }
  \end{center}

  \begin{flushright}
    \tiny{\textbf{source:} Thailand Marine Department (http://www.md.go.th/md/) }
  \end{flushright}

  \note{
    \noindent\textbf{Purpose:} infrastructure sea port in Thailand\par
    \noindent\textbf{Leading Question:} What would you measure before selecting a Thai seaport for a product?\par
    \noindent\textbf{Teaching cue:} Compare inland access, service frequency, handling capability, and cost.\par
    \begin{itemize}
      \item \textbf{Key Gov Ports:} Bangkok (Dry) + Lamchabang (Dry) + Maptapute (CHE)
      \item \textbf{Key Private Port:} SSI, Seeracha
    \end{itemize}
  }
  \end{frame}

%==============================================================================================================
\begin{frame}
  \frametitle{Package Carriers: eCommerce backbone}

  \begin{itemize}
    \item \textbf{Important:} light ($\leq$ 60kg), \alert{consolidation}, multiple-mode
    \item \textbf{Pricing:} weight, class, destination
    \item \textbf{Opportunity:} \alert{responsive}
    \item \textbf{Issues:} \alert{transit facility}
    \item \textbf{Carriers:} UPS, FedEx, DHL, USPS, Kerry, Flash, J\&C
  \end{itemize}

  \note{
    \noindent\textbf{Purpose:} Overview package\par
    \noindent\textbf{Leading Question:} Which elements of a parcel network matter for reliable e-commerce delivery?\par
    \noindent\textbf{Teaching cue:} Trace sorting, linehaul, delivery, and returns using a timing or service measure.\par
    \begin{itemize}
      \item \textbf{Important to eBusiness} risk of super-market place; cashvan
    \end{itemize}
  }
  \end{frame}

%==============================================================================================================
\begin{frame}
  \frametitle{Other Systems}

  \textbf{Pipeline:}
  \begin{itemize}
    \item \textbf{Important:} dominated by crude petroleum, refined petroleum products, natural gas
    \item \textbf{Issues:} \alert{security}, environmental, construction cost, inventory
  \end{itemize}

  \NextHalfLine
  \textbf{Intermodal:}
  \begin{itemize}
    \item \textbf{Important:} multiple-mode to avoid bottleneck in network
    \item \textbf{Combination:} Train-Truck, Water-Train-Truck
    \item \textbf{Opportunity:} high profit margin, unitization
    \item \textbf{Issues:} \alert{communication}, Land bride and Containerization
   \end{itemize}


  \note{
    \noindent\textbf{Purpose:} Overview of other mode\par
    \noindent\textbf{Leading Question:} Which unusual shipment would justify a specialized mode or service?\par
    \noindent\textbf{Teaching cue:} Specify the requirement that a standard mode cannot meet and compare alternatives.\par
    \begin{itemize}
      \item Biggest disadvantage is that carriers are reluctant to participate.
      \item Cultural bias towards using only one mode and this makes change more difficult
    \end{itemize}
  }
  \end{frame}


%%%==============================================================================================================
%%\begin{frame}
%%  \frametitle{Ranking of transportation Modes}
%%
%%  \begin{center}
%%  \begin{tabular}{lccccc} \hline
%%    Size  & Lot       & Safety    & In-Transit  & Transportation  & Transportation \\
%%          & Inventory & Stock     & Cost        & Time            & Cost        \\  \hline
%%    Rail  & 5         & 5         & 5           & 2               & 5 \\
%%    Truckload& 4      & 4         & 4           & 3               & 3 \\
%%    LTL   & 3         & 3         & 3           & 4               & 4 \\
%%    Package& 1        & 1         & 1           & 6               & 1 \\
%%    Air   & 2         & 2         & 2           & 5               & 2 \\
%%    Water  & 6        & 6         & 6           & 1               & 6 \\    \hline
%%  \end{tabular}
%%  \end{center}
%%
%%  \note{
%%    \begin{itemize}
%%      \item
%%    \end{itemize}
%%  }
%%  \end{frame}
%
%==============================================================================================================
\section[Issue]{Current Issues on Transportation}
%==============================================================================================================
\begin{frame}
  \frametitle{Mode Comparison: Historical Illustration}

  \begin{center}
  \footnotesize{
  \begin{tabular}{lrrrrr} \hline
                   & Road                 & Rail        & Air         & Water           & Pipeline        \\  \hline
    \$/ton-mile    & medium               & low/med     & high        & \alert{low/v. low}      & v. low  \\
    Speed (mph)    & 0-60                 & 0-50        & \alert{0-600}& 0-20           & 0-5 \\
    Frequency      & \alert{v.good}               & good        & good        & limited         & \alert{continuous} \\
    Network        & \alert{extensive}    & limited     & limited     & \alert{restricted}      & dedicated \\
    Security       & good                 & average     & average+    & poor          & good \\
    Reliability    & v.good               & good        & \alert{v.good}      & limited     & v. good \\
    Advantage      & mixed/\# vehicle     & heavy load  & low inv.    & cost    & continuous \\
    Limitation     & $\bullet$ loading    & $\bullet$ connectivity    & $\bullet$ cost     & $\bullet$ min qty    & $\bullet$ liquid\\
                   & $\bullet$ dimension  & $\bullet$ network         & $\bullet$ weather  & $\bullet$ speed      & $\bullet$ capital \\
                   & $\bullet$ capacity   &                           & $\bullet$ capacity & $\bullet$ seasonality&   \\ \hline
  \end{tabular}
  }
  \end{center}

  \begin{flushright}
    \tiny{\textbf{source:} Adopted from \textsl{Frazelle} 2001. }
  \end{flushright}

  \note{
    \noindent\textbf{Purpose:} Summary of all mode\par
    \noindent\textbf{Leading Question:} Are the tabulated speeds and costs enough to choose a mode?\par
    \noindent\textbf{Teaching cue:} Treat ranges as illustrative and compare connecting time, service reliability, and a real rate quote.\par
    \begin{itemize}
      \item \textbf{Recall} Limitation \& red text
    \end{itemize}
  }
  \end{frame}


%==============================================================================================================
\begin{frame}
  \frametitle{Article: The Threat of Global GridLock \forIE}

    \begin{itemize}
      \item \textbf{Purpose:} warning about transportation infrastructure
      \item \textbf{Infrastructure:} truck, rail, ship, air
      \item \textbf{Effects to company:} SC is longer and more complex $\rightarrow$ severe congestions
      \begin{itemize}
        \item delay has cascade effect
        \item distant souring is deterred by  high fuel cost
        \item rapid change product $\rightarrow$ ``double dips'' \NextHalfLine
      \end{itemize}
      \item \textbf{Remedies:} goal is to maximize profit, not minimize costs
      \begin{itemize}
        \item cost \textbf{VS} time
        \item premium service (``hot-hatching'', door-to-door, reserve cap)
        \item source closer, using Air, using best practice
      \end{itemize}
    \end{itemize}

  \note{
    \noindent\textbf{Purpose:} Recent review of modes of transportation\par
    \noindent\textbf{Leading Question:} What evidence would support a congestion claim for a current route?\par
    \noindent\textbf{Teaching cue:} Check the reading date, location, measurement, and alternative routes before using it.\par
    \begin{itemize}
      \item \textbf{hot-hatching} key inventory concept that lead to division
      \item \textbf{door-to-door} key inventory concept that lead to division
      \item \textbf{double dips} key inventory concept that lead to division
    \end{itemize}  }
  \end{frame}

%==============================================================================================================
\section[Thappline]{Pipeline Infrastructure in Thailand: a case study}
%=================================================================================================
\begin{frame}
  \frametitle{Pipeline Network in Thailand \forIE}

  \begin{tabular}{cc}
    \includegraphics [width=0.45\textwidth]{\figureLocation{scm_mode/existingPipeline.png}}
    &
    \includegraphics [width=0.45\textwidth]{\figureLocation{scm_mode/thappline.png}}
  \end{tabular}

  \note{
    \noindent\textbf{Purpose:} Overview of pipeline infrastructure in Thailand\par
    \noindent\textbf{Leading Question:} When would a dedicated pipeline justify its fixed investment?\par
    \noindent\textbf{Teaching cue:} Compare steady demand, alternatives, capacity use, and operational constraints at a high level.\par
    \begin{itemize}
      \item \textbf{Pipeline:} limited on ground for refined oils and gas transportation; extensive in water for gas and electricity production
      \item \textbf{Thappline:} Thai Petrolumn pipeline found by TH gov and PTT. Currently, it privatizes and operates ground pipeline.
    \end{itemize}  }
  \end{frame}

%=================================================================================================
\begin{frame}
  \frametitle{Pipeline Network in Thailand}

  \begin{itemize}
    \item \textbf{Background:} Thappline (1990; former subsidiary of PTT)
    \item \textbf{Products} mixed product (GB-1, GB-2, ULG, \& HSD) in single pipe
    \item \textbf{Piping Tech:} 24"-12" diameter; 76 Block Valve Station
    \begin{itemize}
      \item \textsl{Open Cut:} removing soil surface 1.5m before place pipe for general
      \item \textsl{Bored Crossing/Pipe jacking:} pushing pipeline underground for small intersection
      \item \textsl{Horz Drilling:} several drilling deep with guide sys for river
    \end{itemize}
    \item \textbf{Piping Route:} existing major road and railway
  \end{itemize}

  \note{
    \noindent\textbf{Purpose:} Overview Inventory in SCM\par
    \noindent\textbf{Leading Question:} Which network trade-off does the Thappline case illustrate?\par
    \noindent\textbf{Teaching cue:} Discuss right-of-way and throughput qualitatively, then check the dated details against the case source.\par
    \begin{itemize}
      \item \textbf{Operation:} Given time, each refinery must pump specific type of oil in pipeline with meter $\rightarrow$ sale volume
      \item \textbf{Pig} equipment to separate and track oil
      \item \textbf{Oil Type} GB-1: Gasohol Benzene 91/ GB-2: 95 ULG: Unleaded gasolene / HSD: High speed Diesel
    \end{itemize}
  }
  \end{frame}

%================================================================================================-

\begin{frame}
  \frametitle{Reference}

  \begin{footnotesize}
    \bibliography{\refLocation{Coordination}}
    %\bibliography{\refLocation{logisticProvider}}
  \end{footnotesize}

  \note{
    \noindent\textbf{Purpose:} Show reference\par
    \noindent\textbf{Leading Question:} What would you check before reusing a figure, fact, or model from this lecture?\par
    \noindent\textbf{Teaching cue:} Trace a citation to its original source and verify the year, scope, and assumptions.\par
    \begin{itemize}
      \item Confirm the source and edition before assigning a reading.
    \end{itemize}
  }

\end{frame}

%=================================================================================================
\end{document}



%%=================================================================================================
\begin{frame}
  \frametitle{Proposed Oil Pipeline \& Stations}

  \small{
    \begin{tabular}{cl}
      %\multirow{2}{*}{
        \includegraphics[width=0.40\textwidth]{\figureLocation{pipelineProject.png}}
      %}
    &
    \includegraphics [width=0.45\textwidth]{\figureLocation{oilStationLayout.png}} \\
    & \textbf{Center:} RA, SA, WA \\
    & \textbf{North:} LA, KA\\
    & \textbf{Northeast:} NA, KH
  \end{tabular}
  }

  \note{
    \begin{itemize}
      \item Pipeline of ligure 67.0 km
      \item Pipeline of gas 350.0 km	
    \end{itemize}
  }
  \end{frame}

%==============================================================================================================
\section[LSP]{Logistics Service Provider}
%=================================================================================================
\begin{frame}
  \frametitle{Why Outsourcing and Types of LSP}

  \begin{itemize}
    \item \textbf{Reasons:} reduce costs, gain expertise | flexibility  \& core competency
    \item \textbf{Other consideration:} computability, politic \& suitability (FoodStar)
    \item \textbf{Concerns:} loss control, no information \& knowledge, \alert{partnership}
    \item \textbf{Types:}
    \begin{itemize}
      \item \textsc{asset-based:} physical asset, e.g., refrigerated trucks; chemical warehouse (WHA)
      \item \textsc{management-based:} knowledge, connection, IT
      \item \textsc{integration-based:} single service providers,
      \item \textsc{administrative-based:} e.g. payment,
    \end{itemize}
  \end{itemize}
  \begin{flushright}
    \tiny{\textbf{source:} \textsl{Bourlakis, M. and Melewar, T.} 2011.~\cite{bourlakis2011marketing}~~ }
  \end{flushright}


  \note{
    \begin{itemize}
      \item
    \end{itemize}
  }
  \end{frame}

%%=================================================================================================
\begin{frame}
  \frametitle{Survey of LSP usage in W.Europe \& US}

  \begin{center}
    \includegraphics[width=0.35\textwidth]{\figureLocation{3plWE_US.png}} \\
    \includegraphics[width=0.80\textwidth]{\figureLocation{3plSatis.png}}
  \end{center}

  \begin{flushright}
    \tiny{\textbf{source:} \textsl{Lieb, R, et al.} 1993.~\cite{lieb1993third}~~~ }
  \end{flushright}

  \note{
    \begin{itemize}
      \item
    \end{itemize}
  }
  \end{frame}

%%=================================================================================================
\begin{frame}
  \frametitle{Managing Expectation}

  \footnotesize{
  \begin{tabular}{lllll} \hline \hline
                   & \textbf{Expectation} $\Rightarrow$ & \textbf{Realization} $\Rightarrow$ & \textbf{Stagnation} $\Rightarrow$    & \textbf{Frustration}  \\ \cline{2-5}
    \textbf{State} & $\bullet$ excitement               & $\bullet$ diff. imple.        & $\bullet$ operating         & $\bullet$ separated \\
                   & $\bullet$ enthusiasm               &                               & $\bullet$ not expected      & \\[6pt]  \hline
    \textbf{Owner} & $\bullet$ expect turnover          & $\bullet$ cost$+ \rightarrow$ ROI & $\bullet$ do kaizen     & $\bullet$ need change \\
                   &                                    & $\bullet$ feel distrust/bad   & $\bullet$ little improvement& $\bullet$ find new provider \\[6pt] \hline
    \textbf{Provider}& $\bullet$ put 'best' team        & $\bullet$ cost$+ \rightarrow$ time& $\bullet$ commun. brkDwn& $\bullet$ know oper.  \\
                   &                                    &                               & $\bullet$ not enough kaizen & $\bullet$ hi bidding                 \\ \hline    \hline
  \end{tabular}
  }

  \NextHalfLine
  \textsc{Warnings:} outsourcing $\approx$ marriage
  \begin{itemize}
    \item understand \alert{current operation}, never outsource problems
    \item understand \alert{cost structure} before/after outsource
    \item understand \alert{service provider market}, consult/search for 'right' LSP
  \end{itemize}


  \note{
      \begin{itemize}
        \item \textbf{Dimension:} Fiance, Long-term relationship, Operation Performance, Risk
        \item \textbf{distances:} road distance (know exact location), point-point
        \item \textbf{service levels:} specific safety stock
        \item \textbf{Facility cost:} time horizon, discounted factor
      \end{itemize}
   }
 \end{frame}

%%=================================================================================================
\begin{frame}
  \frametitle{LSP selection}

  \begin{itemize}
    \item \textbf{Determinant:} \alert{Compatibility}, Cost, Quality, Reputation
    \item \textbf{Specifically,} performance, data sharing, assets, experience, market shares, range of services, network
    \item \textbf{How to find LSP?}
    \begin{itemize}
      \item[1] determine service objectives $\rightarrow$ gap \& specification
      \item[2] list of potential LSPs
      \item[3] develop and evaluate RFI $\rightarrow$ RFQ
      \item[4] survey and site visits $\rightarrow$ sign contract
      \item[5] evaluate performance
    \end{itemize}
  \end{itemize}

  \NextHalfLine
  \begin{center}
  \textsc{Decision Matrix}\\[5pt]
  \footnotesize{
  \begin{tabular}{llrr} \hline
                                  &~& \multicolumn{2}{c}{Own Expertise} \\    \cline{3-4}
                                  &~& \textbf{Low} & \textbf{High} \\ \hline
    \multirow{2}{*}{Value} & \textbf{High} &  partnership outsourcing & do not outsource     \\
                           & \textbf{Low}  &  transactional outsourcing & depends on finance    \\ \hline
  \end{tabular}
  \NextHalfLine
  \begin{itemize}
    \item \textbf{Partnership:} long term \& focus on efficiency (PharmInd \& BioTech)
    \item \textbf{Transactional/Tactical:} short term \& focus on cost  (PharmaInd--1970)
  \end{itemize}
  }
  \end{center}


  \begin{flushright}
    \tiny{\textbf{source:} \textsl{Vitasek, T.} 2010.~\cite{vitasek2010marketing}~~ }
  \end{flushright}


  \note{
      \begin{itemize}
        \item \textbf{Dimension:} Fiance, Long-term relationship, Operation Performance, Risk
        \item \textbf{distances:} road distance (know exact location), point-point
        \item \textbf{service levels:} specific safety stock
        \item \textbf{Facility cost:} time horizon, discounted factor
    \item \textbf{Partnership:} long term \& efficiency, share culture, technology, infrastructure (PharmInd \& BioTech)
    \item \textbf{Transactional/Tactical:} short term| cost focus, no collaborative nor business alignment (PharmaInd--1970)
      \end{itemize}
   }
 \end{frame}

%%=================================================================================================
%\begin{frame}
%  \frametitle{Tailored Network}
%
%  \begin{itemize}
%    \item \textbf{What:} customized network
%    \item \textbf{Goal:} using appropriate options in each situation
%    \item \textbf{Challenge:} no concrete network, required near-perfect planning and execution
%    \item \textbf{Criteria:} customer density and size, distance, product demand, value
%  \end{itemize}
%
%  \note{
%      \begin{itemize}
%        \item A company focuses on inbound b/c easy to pressure supplier than customers
%      \end{itemize}
% }
% \end{frame}
%
%%==============================================================================================================
%\begin{frame}
%  \frametitle{Land Transportation Options}
%
%  \begin{center}
%  \small{
%  \begin{tabular}{rlll} \hline %\hline
%                   & Short distance                  & Medium distance         & Long distance        \\  \hline %\hline
%    High density   & $\bullet$ Private fleet            & $\bullet$ Cross-dock       & $\bullet$ Cross-dock    \\
%                   & with milk runs                  & with milk runs          & with milk runs       \\  %\hline
%    Medium density & $\bullet$ 3PL with milk runs       & $\bullet$ LTL carrier      & $\bullet$ LTL carrier   \\
%                   &                                 &                         &   or package carrier \\  %\hline
%    Low density    & $\bullet$ 3PL with milk runs    & $\bullet$ LTL carrier      & $\bullet$ Package carrier \\
%                   & or LTL carrier                  & or package carrier      & \\    \hline %\hline
%  \end{tabular}
%  }
%  \end{center}
%
%  \note{
%    \begin{itemize}
%      \item
%    \end{itemize}
%  }
%  \end{frame}


%%=================================================================================================
%\section[Modeling]{Distribution Network Modeling in Operation Management}
%\begin{frame}
%  \frametitle{Why do we need a model?}
%  \begin{itemize}
%    \item understand and analyze network
%    \item predict and determine decisions
%  \end{itemize}
%
%  \NextLine
%  \textbf{Related Costs}
%  \begin{itemize}
%    \item \textbf{Production/Purchasing cost:} qualities discount
%    \item \textbf{Transportation cost:} \alert{distance}, weight, volume, \# shipments
%    \item \textbf{Inventory cost:} \alert{service levels}, safety stock, pipeline,
%    \item \textbf{Facility cost:} rental, equipments (e.g. sorting machine, forklifts)
%  \end{itemize}
%
%  \note{
%      \begin{itemize}
%        \item \textbf{so far:} introduction on terminology
%        \item \textbf{distances:} road distance (know exact location), point-point
%        \item \textbf{service levels:} specific safety stock
%        \item \textbf{Facility cost:} time horizon, discounted factor
%      \end{itemize}
%   }
% \end{frame}

%========================================================================================================
\begin{frame}
  \frametitle{Key Questions in Distribution Design}

  \begin{itemize}
    \item \textbf{Strategic level}
    \begin{itemize}
      \item<2-> Mode of transportation, In-house logistics
      \item<2-> \alert{Locations and numbers of transit facilities}
      \item<2-> Sizes of transit facilities and their equipments \NextLine
    \end{itemize}
    \item \textbf{Operation levels}
    \begin{itemize}
      \item<3-> Storage space of each product and its inventory policy
      \item<3-> Routing and time window of each destination
      \item<3-> Fleet management (crew/driver, empty truck)
    \end{itemize}
  \end{itemize}

  \note{
    \begin{itemize}
      \item
    \end{itemize}
  }
\end{frame}

%========================================================================================================
  \begin{frame}
  \frametitle{Distribution Network Design}

  \textbf{Components}
  \begin{itemize}
    \item \textbf{Owner:} what factors are his \alert{total transportation costs}?
    \item \textbf{Nodes:} location (customers, supplier), flows (weight, volume), service level
    \item \textbf{Candidate:}
    \begin{itemize}
      \item given list $\rightarrow$ integer problem
      \item to-be-determined $\rightarrow$ Weber's problem
    \end{itemize}
    \item \textbf{transportation cost:} fixed costs, variable costs
  \end{itemize}

    \NextLine
    \textbf{Examples}
    \begin{itemize}
      \item single facility (1D, 2D)
      \item multiple facility
    \end{itemize}


  \note{
    \begin{itemize}
      \item
    \end{itemize}
  }
\end{frame}


Twenty-foot Equivalent Unit (TEU) smallest standard shipping container.                                                        	

Demurrage
DWT (Dead Weight tonnage)
short ton/ long ton/ matric ton (tonne)


